\documentclass[11pt,letterpaper]{article}
\usepackage[utf8]{inputenc}
\usepackage[T1]{fontenc}
\usepackage[spanish,es-tabla,es-nodecimaldot,es-noquoting,es-noshorthands]{babel}
\usepackage[top=2.3cm,bottom=2.2cm,left=2.1cm,right=2.1cm]{geometry}
\usepackage{graphicx,booktabs,tabularx,array,xcolor,tikz,enumitem,fancyhdr,hyperref,amsmath,amssymb}
\usetikzlibrary{arrows.meta,positioning,fit,backgrounds,shapes.geometric,babel}

% Paleta institucional UPV
\definecolor{upvmorado}{HTML}{5B146B}
\definecolor{upvazul}{HTML}{17384D}
\definecolor{upvverde}{HTML}{008D79}
\definecolor{upvnaranja}{HTML}{E67817}
\definecolor{grisclaro}{HTML}{F8FAFC}
\definecolor{grisborde}{HTML}{CBD5E1}
\definecolor{gristexto}{HTML}{475569}
\definecolor{azulsuave}{HTML}{EFF6FF}
\definecolor{moradosuave}{HTML}{FAF5FF}

\hypersetup{
    colorlinks=true,
    urlcolor=upvverde,
    linkcolor=upvazul,
    pdftitle={Proyecto final: Sistema inteligente para la generación de horarios de ITIID - Manual de Usuario},
    pdfauthor={Equipo ITIID - Universidad Politécnica de Victoria}
}

\setlength{\parindent}{0pt}
\setlength{\parskip}{5pt}
\emergencystretch=1.8em
\setlist[itemize]{itemsep=2pt,topsep=3pt,leftmargin=*}
\setlist[enumerate]{itemsep=3pt,topsep=3pt,leftmargin=*}
\renewcommand{\arraystretch}{1.22}
\addto\captionsspanish{\renewcommand{\contentsname}{Índice General}}

% Encabezados y pies de página
\pagestyle{fancy}
\fancyhf{}
\fancyhead[L]{%
    \raisebox{-0.18\height}{\includegraphics[height=0.52cm]{img/logo_upv_morado.png}}\hspace{0.22cm}%
    {\footnotesize\bfseries\color{upvmorado} UPV $\cdot$ ITIID}\hspace{0.15cm}{\footnotesize\color{gristexto}| Sistema Inteligente de Horarios}%
}
\fancyhead[R]{\footnotesize\color{gristexto} Manual de Usuario y Extensión}
\fancyfoot[L]{\footnotesize\color{gristexto} Dr. José Fidencio López Luna $\cdot$ Fundamentos de IA}
\fancyfoot[R]{\footnotesize\bfseries\color{upvazul} Página \thepage}
\renewcommand{\headrulewidth}{0.5pt}
\renewcommand{\footrulewidth}{0.3pt}
\setlength{\headheight}{20pt}

% Comandos de formato y cajas
\newcommand{\seccionUPV}[1]{%
    \section{\texorpdfstring{\textcolor{upvmorado}{#1}}{#1}}%
    \vspace{-0.22cm}{\color{upvnaranja}\rule{\textwidth}{1.1pt}}\vspace{0.12cm}%
}

\newcommand{\subseccionUPV}[1]{%
    \subsection{\texorpdfstring{\textcolor{upvazul}{#1}}{#1}}%
}

\newcommand{\cajaNota}[2]{%
    \par\medskip\noindent
    \begin{tikzpicture}
        \node[
            fill=moradosuave,
            draw=upvmorado!35,
            rounded corners=3pt,
            inner xsep=11pt,
            inner ysep=8pt,
            text width=\dimexpr\textwidth-24pt\relax,
            align=justify
        ] (box) {\textbf{\color{upvmorado}#1:} #2};
        \draw[upvmorado, line width=3pt, rounded corners=1pt] (box.north west) -- (box.south west);
    \end{tikzpicture}
    \par\medskip
}

\newcommand{\cajaCodigo}[1]{%
    \par\smallskip\noindent
    \colorbox{grisclaro}{%
        \parbox{\dimexpr\linewidth-2\fboxsep\relax}{%
            \vspace{3pt}{\ttfamily\small #1}\vspace{3pt}%
        }%
    }\par\smallskip
}

\newcounter{figmanual}
\newcommand{\captura}[3][0.88]{%
    \par\medskip\noindent
    \begin{minipage}{\textwidth}
        \centering
        \setlength{\fboxsep}{2pt}%
        \setlength{\fboxrule}{0.6pt}%
        \fcolorbox{grisborde}{white}{\includegraphics[width=#1\textwidth]{img/#2}}\\[4pt]
        \refstepcounter{figmanual}%
        {\footnotesize\textbf{\color{upvazul}Figura \thefigmanual.} {\color{gristexto}#3}}
    \end{minipage}
    \par\medskip
}

\begin{document}

% ==============================================================================
% PORTADA INSTITUCIONAL UPV
% ==============================================================================
\begin{titlepage}
    \centering
    \vspace*{0.1cm}
    \includegraphics[height=2.35cm]{img/logos_upv_bis.png}\\[0.35cm]
    {\color{upvmorado}\rule{\textwidth}{2pt}}\\[-0.22cm]
    {\color{upvnaranja}\rule{\textwidth}{0.9pt}}\\[0.35cm]

    {\Large\bfseries\color{upvazul} UNIVERSIDAD POLITÉCNICA DE VICTORIA\par}
    \vspace{0.12cm}
    {\normalsize\bfseries\color{gristexto} Ingeniería en Tecnologías de la Información e Innovación Digital (ITIID)\par}
    \vspace{0.35cm}

    \includegraphics[height=2.65cm]{img/logo_upv_morado.png}\\[0.35cm]

    {\LARGE\bfseries\color{upvmorado} Proyecto final: Sistema inteligente para la\\generación de horarios de ITIID\par}
    \vspace{0.28cm}
    {\large\bfseries\color{upvazul} Manual de Usuario, Guía de Descarga e Instalación\\y Manual Técnico de Modificación y Extensión\par}
    \vspace{0.45cm}

    \begin{tikzpicture}
        \node[
            fill=grisclaro,
            draw=grisborde,
            rounded corners=4pt,
            inner xsep=14pt,
            inner ysep=9pt,
            text width=0.88\textwidth
        ] {%
            \begin{tabularx}{\linewidth}{@{}lX@{}}
                \textbf{\color{upvazul}Asignatura:} & Fundamentos de Inteligencia Artificial (7.\textsuperscript{o} Cuatrimestre) \\
                \textbf{\color{upvazul}Dr. encargado:} & \textbf{Dr. José Fidencio López Luna} \\
                \textbf{\color{upvazul}Institución:} & Universidad Politécnica de Victoria (UPV) --- Modelo BIS
            \end{tabularx}
        };
    \end{tikzpicture}

    \vspace{0.45cm}
    {\normalsize\bfseries\color{upvmorado} INTEGRANTES DEL EQUIPO\par}
    \vspace{0.18cm}

    \begin{tikzpicture}
        \node[
            fill=white,
            draw=upvmorado!35,
            rounded corners=4pt,
            inner xsep=16pt,
            inner ysep=10pt,
            text width=0.82\textwidth,
            align=center
        ] {%
            \centering
            \begin{tabular}{@{}c@{}}
                Axel Librado Aguilera Alonso \\
                Germán David Zavala Castro \\
                Vasti Norely Ruiz Báez \\
                Yahir Alejandro Roque Martínez \\
                José Guadalupe Cornejo Alva \\
                David Alejandro Ibarra Castañeda \\
                Jesús Alberto Tovar Ramírez
            \end{tabular}
        };
    \end{tikzpicture}

    \vfill
    {\color{upvmorado}\rule{0.55\textwidth}{0.6pt}}\\[0.15cm]
    {\small\color{gristexto} Ciudad Victoria, Tamaulipas, México $\cdot$ Noviembre de 2026\par}
\end{titlepage}

% ==============================================================================
% ÍNDICE GENERAL
% ==============================================================================
\pagenumbering{roman}
\tableofcontents
\newpage
\pagenumbering{arabic}

% ==============================================================================
% 1. INTRODUCCIÓN Y ARQUITECTURA DEL SISTEMA
% ==============================================================================
\seccionUPV{Introducción y Arquitectura del Sistema}

El \textbf{Sistema inteligente para la generación de horarios de ITIID} es una plataforma web híbrida desarrollada para automatizar la planeación, validación, comparación experimental y reprogramación de horarios académicos del programa de \textbf{Ingeniería en Tecnologías de la Información e Innovación Digital (ITIID)} en la \textbf{Universidad Politécnica de Victoria (UPV)}.

La plataforma conserva intactos los tres módulos base de administración y resolución por \textit{backtracking} en C++ (\textbf{Inicio}, \textbf{Datos} y \textbf{Horarios}) e incorpora una suite completa de \textbf{Inteligencia Artificial} compuesta por tres módulos especializados: \textbf{Optimizador IA}, \textbf{Asistente LLM} y \textbf{Comparativa y Ética}.

\cajaNota{Regla central de diseño (IA Híbrida Segura)}{Las asignaciones de clases nunca quedan a criterio libre de un modelo de lenguaje. Los horarios son construidos exclusivamente por motores formales de optimización (\textbf{OR-Tools CP-SAT}, \textbf{Algoritmo Genético}, \textbf{Búsqueda Local} o el solver clásico en \textbf{C++}) y auditados por un \textbf{Validador Común} determinista. El \textbf{LLM} actúa únicamente como traductor de preferencias en español con revisión humana obligatoria (\textit{human-in-the-loop}), explicador de conflictos y traductor de solicitudes de cambio.}

\subseccionUPV{Arquitectura General de Componentes}

El sistema opera de manera 100\% local mediante tres capas desacopladas que se comunican por HTTP/JSON y tuberías estándar (\texttt{stdin} y \texttt{stdout}):

\begin{center}
\begin{tikzpicture}[
    font=\scriptsize,
    box/.style={draw=upvazul, rounded corners=3pt, align=center, fill=azulsuave, minimum height=1.0cm, inner sep=5pt},
    iabox/.style={draw=upvmorado, rounded corners=3pt, align=center, fill=moradosuave, minimum height=1.0cm, inner sep=5pt},
    extbox/.style={draw=upvverde, rounded corners=3pt, align=center, fill=green!5, minimum height=0.9cm, inner sep=5pt},
    arr/.style={-{Latex[length=2mm]}, thick, draw=upvazul}
]
    \node[box, text width=4.3cm] (fe) at (0,0) {\textbf{Frontend Web (React 19 + Vite)}\\Puerto \texttt{localhost:5173}\\Tailwind CSS $\cdot$ Recharts $\cdot$ LocalStorage};
    \node[box, text width=4.6cm] (api) at (5.8,0) {\textbf{Servidor Puente (\mbox{Node.js} + Express)}\\Puerto \texttt{localhost:3001}\\\texttt{/api/solve} y \texttt{/api/ia/action}};
    \node[box, text width=4.3cm] (cpp) at (11.6,1.55) {\textbf{Motor Clásico C++}\\\texttt{backend/solver.exe}\\Backtracking + Poda horizontal};
    \node[iabox, text width=4.3cm] (py) at (11.6,-0.65) {\textbf{Motor IA Python (\texttt{cli.py})}\\CP-SAT $\cdot$ Genético $\cdot$ Búsq. Local\\Validador Común $\cdot$ Excel UPV};
    \node[extbox, text width=4.3cm] (llm) at (11.6,-3.15) {\textbf{Servicio LLM 100\% Local}\\Ollama Local (\texttt{127.0.0.1:11434})\\Modelo \texttt{qwen2.5:3b} (Sin APIs externas)};

    \draw[arr, <->] (fe) -- node[above]{HTTP JSON} (api);
    \draw[arr, <->] (api.east) -- ++(0.45,0) |- (cpp.west);
    \draw[arr, <->] (api.east) -- ++(0.45,0) |- (py.west);
    \draw[arr, <->] (py) -- (llm);
\end{tikzpicture}
\end{center}

\begin{itemize}
    \item \textbf{Frontend (\texttt{frontend/}):} Aplicación SPA en React 19 con Vite 7, estilos en Tailwind CSS, gráficas estadísticas en Recharts, visualización topológica con \texttt{react-force-graph-2d} y exportación cliente en PDF/Excel (\texttt{jspdf}, \texttt{xlsx}).
    \item \textbf{Servidor Puente (\texttt{server/server.js}):} API REST en Express (puerto \texttt{3001}) que recibe peticiones del navegador, ejecuta procesos hijos hacia C++ o Python con control de concurrencia y tiempos límite, e incluye respaldo automático hacia Python CP-SAT si faltan librerías de C++ en el equipo.
    \item \textbf{Motor Clásico (\texttt{backend/solver.cpp}):} Algoritmo exacto de búsqueda con retroceso (\textit{backtracking}) y distribución horizontal de sesiones semanales.
    \item \textbf{Motor de Inteligencia Artificial (\texttt{backend\_ia/}):} Paquete modular en Python que integra el modelo exacto \textbf{Google OR-Tools CP-SAT} (R2), \textbf{Algoritmo Genético} multicorrida (R3), \textbf{Búsqueda Local MRV + Hill Climbing} (R4), \textbf{Validador Oficial} compartido, generador de \textbf{Sábana Oficial UPV en Excel} (\texttt{openpyxl}, R1) y el servicio de lenguaje natural con anonimización automática (R5/R6).
\end{itemize}

% ==============================================================================
% 2. QUÉ SE TIENE QUE DESCARGAR (REQUISITOS DE SOFTWARE Y HARDWARE)
% ==============================================================================
\seccionUPV{Qué Tienes que Descargar (Requisitos Previos)}

Para instalar y ejecutar el proyecto desde cero en cualquier computadora con Windows 10/11 (o Linux/macOS), es necesario descargar las herramientas listadas en la Tabla~\ref{tab:descargas}.

\begin{table}[htbp]
\centering
\small
\caption{Programas y librerías que se deben descargar para ejecutar el sistema.}
\label{tab:descargas}
\begin{tabularx}{\textwidth}{@{}>{\raggedright\arraybackslash}p{3.2cm}>{\raggedright\arraybackslash}p{2.2cm}>{\raggedright\arraybackslash}X>{\raggedright\arraybackslash}p{4.3cm}@{}}
\toprule
\textbf{Componente} & \textbf{Carácter} & \textbf{Para qué sirve en el proyecto} & \textbf{Dónde descargarlo} \\
\midrule
\textbf{1. Node.js LTS} (v18 o superior; rec. v20/v22) & \textbf{Obligatorio} & Ejecuta el servidor Express (\texttt{server/}) e instala/sirve la interfaz React con \texttt{npm} (\texttt{frontend/}). & \url{https://nodejs.org} (instalador \texttt{.msi} LTS) \\
\addlinespace[3pt]
\textbf{2. Python} (v3.11 o v3.12) & \textbf{Obligatorio} & Ejecuta todo el motor \texttt{backend\_ia/} (CP-SAT, Genético, Búsqueda Local, Validador y Excel). & \url{https://www.python.org/downloads/} \\
\addlinespace[3pt]
\textbf{3. Paquetes \texttt{pip}} (\texttt{ortools}, \texttt{openpyxl}, \texttt{pandas}, \texttt{requests}) & \textbf{Obligatorio} & Librerías matemáticas de optimización exacta, lectura/escritura de plantillas Excel UPV y cliente HTTP. & Desde terminal con \texttt{pip install} (ver comando abajo) \\
\addlinespace[3pt]
\textbf{4. Ollama + Modelo \texttt{qwen2.5:3b}} ($\sim$1.9 GB) & \textbf{Obligatorio para LLM} & Ejecuta el modelo de lenguaje 100\% local en GPU/CPU en \texttt{127.0.0.1:11434}, sin claves API ni servicios en la nube. & \url{https://ollama.com/download} o \texttt{winget install Ollama.Ollama} \\
\addlinespace[3pt]
\textbf{5. Compilador C++ \texttt{g++}} (MinGW-w64) & \textbf{Opcional} & Solo se requiere si deseas modificar y recompilar \texttt{backend/solver.cpp}. El binario \texttt{solver.exe} ya viene incluido. & \url{https://winlibs.com} o MSYS2 \\
\addlinespace[3pt]
\textbf{6. Navegador moderno} (Chrome / Edge) & \textbf{Obligatorio} & Renderiza la interfaz en \texttt{http://localhost:5173}, el grafo 2D en Canvas y las descargas PDF/Excel. & Google Chrome o Microsoft Edge \\
\bottomrule
\end{tabularx}
\end{table}

\subseccionUPV{Requisitos Mínimos y Recomendados de Hardware}
\begin{itemize}
    \item \textbf{Procesador:} Intel Core i3 / AMD Ryzen 3 de 4 núcleos (Core i5 / Ryzen 5 o superior recomendado para comparativas multicorrida y Ollama local).
    \item \textbf{Memoria RAM:} 4 GB mínimo para interfaz web + motores de optimización; \textbf{8 GB o más} si se ejecuta el modelo local \texttt{qwen2.5:3b} en Ollama.
    \item \textbf{Espacio libre en disco:} 600 MB para el proyecto y dependencias (\texttt{node\_modules} + paquetes Python), más 2.2 GB adicionales si se descarga el modelo local de Ollama.
    \item \textbf{Pantalla:} Resolución mínima \texttt{1366x768} (\texttt{1920x1080} recomendada para apreciar la Sábana General UPV completa).
\end{itemize}

\cajaNota{Importante al instalar Python en Windows}{En la primera pantalla del instalador oficial de Python (\texttt{python-3.12.x-amd64.exe}), marca obligatoriamente la casilla inferior \textbf{``Add python.exe to PATH''} antes de pulsar \textit{Install Now}. Esto permite que Node.js invoque el comando \texttt{python} automáticamente.}

% ==============================================================================
% 3. INSTALACIÓN Y QUÉ HACER PARA QUE FUNCIONE PASO A PASO
% ==============================================================================
\seccionUPV{Instalación y Puesta en Marcha Paso a Paso}

Una vez descargados e instalados \textbf{Node.js} y \textbf{Python}, sigue estos pasos exactos en una terminal (\textit{PowerShell} o \textit{Símbolo del sistema / CMD}) para dejar la aplicación funcionando.

\subseccionUPV{Paso 1: Verificar Node.js y Python en la Terminal}
Abre una terminal nueva y comprueba que los tres comandos respondan con su número de versión:
\cajaCodigo{%
node --version\\
npm --version\\
python --version%
}

\subseccionUPV{Paso 2: Descomprimir el Proyecto e Instalar Librerías de Python}
Extrae la carpeta del proyecto en una ruta accesible (por ejemplo, en el Escritorio) e instala las cuatro librerías del motor de Inteligencia Artificial con un solo comando:
\cajaCodigo{%
pip install ortools openpyxl pandas requests%
}

\subseccionUPV{Paso 3: Instalar Dependencias del Servidor y del Frontend}
Desde la carpeta raíz del proyecto (\texttt{university\_time\_tabbling\_problem\_ed-main}), instala los paquetes de Node.js tanto en \texttt{server/} como en \texttt{frontend/}:
\cajaCodigo{%
cd server\\
npm install\\[4pt]
cd ../frontend\\
npm install\\
cd ..%
}

\subseccionUPV{Paso 4: Instalar Ollama y Descargar el Modelo Local}
El sistema opera \textbf{exclusivamente con Ollama Local} (sin APIs externas ni claves de pago). Instala Ollama y descarga el modelo ligero recomendado:
\cajaCodigo{%
winget install Ollama.Ollama\\
ollama pull qwen2.5:3b%
}
Al iniciar \texttt{server/server.js} o \texttt{run\_project.bat}, el sistema levanta automáticamente el servidor local de Ollama en \texttt{http://127.0.0.1:11434}.

\subseccionUPV{Paso 5: Cómo Iniciar la Aplicación}

Existen dos maneras de encender el sistema cada vez que vayas a utilizarlo:

\textbf{Opción A --- Inicio automático en Windows (1 clic, recomendado):}\\
En el Explorador de Archivos, entra a la carpeta raíz del proyecto y haz \textbf{doble clic sobre \texttt{run\_project.bat}}. El lanzador ejecuta automáticamente cuatro tareas:
\begin{enumerate}
    \item Verifica que exista \texttt{backend\textbackslash solver.exe} (o lo compila si detecta \texttt{g++}).
    \item Detecta si tienes Ollama instalado y levanta \texttt{ollama serve} en segundo plano (\texttt{127.0.0.1:11434}).
    \item Abre la ventana \textbf{Node Server} corriendo \texttt{npm start} en \texttt{http://localhost:3001}.
    \item Abre la ventana \textbf{Frontend} corriendo \texttt{npm run dev} en \texttt{http://localhost:5173} y lanza tu navegador web automáticamente tras 5 segundos.
\end{enumerate}

\textbf{Opción B --- Inicio manual en dos terminales (Windows, Linux o macOS):}\\
Si trabajas desde un editor como VS Code o en Linux/macOS, abre dos terminales en paralelo:

\textit{Terminal 1 --- Servidor Backend y Puente IA (dejar abierta):}
\cajaCodigo{%
cd server\\
npm start%
}
Verás el mensaje: \texttt{Server running on http://localhost:3001}.

\textit{Terminal 2 --- Interfaz Web Vite/React (dejar abierta):}
\cajaCodigo{%
cd frontend\\
npm run dev%
}
Abre tu navegador en \textbf{\texttt{http://localhost:5173}}. Para apagar el sistema al terminar, presiona \texttt{Ctrl + C} en ambas terminales o cierra las ventanas abiertas por \texttt{run\_project.bat}.

\newpage
% ==============================================================================
% 4. MANUAL DE USUARIO: MÓDULOS BASE DEL SISTEMA
% ==============================================================================
\seccionUPV{Manual de Usuario: Módulos Base (Inicio, Datos y Horarios)}

Al ingresar a \texttt{http://localhost:5173}, la interfaz muestra a la izquierda la barra lateral (\textbf{Sidebar}) dividida en dos bloques: los módulos administrativos originales (\textbf{Inicio}, \textbf{Datos}, \textbf{Horarios}) y la suite de \textbf{Inteligencia Artificial} (\textbf{Optimizador IA}, \textbf{Asistente LLM}, \textbf{Comparativa y Ética}).

\subseccionUPV{Módulo Inicio / Vista General (\texttt{/})}

El tablero inicial presenta en tiempo real el estado global de los catálogos guardados en el navegador (\texttt{localStorage}): total de materias activas, profesores registrados y matrícula acumulada, acompañados de una gráfica de barras de alumnos por cuatrimestre y una gráfica de anillo con el reparto de materias del plan ITIID.

\captura{01_dashboard_inicio.png}{Tablero principal (Vista General) con indicadores globales, gráficas por cuatrimestre y menú lateral completo.}

\subseccionUPV{Módulo Gestión de Datos (\texttt{/data})}

Antes de generar un horario personalizado, aquí se administran los tres insumos fundamentales divididos en pestañas superiores:

\textbf{1. Pestaña Profesores (Catálogo Oficial Sep--Dic 2026 y Materias Compartidas):} Incluye precargados los \textbf{32 profesores reales del periodo Sep--Dic 2026} clasificados por categoría (\textbf{FTC} Tiempo Completo, \textbf{PA} Asignatura e \textbf{Inglés}) junto con sus ventanas horarias y materias asignadas. Desde el formulario superior o el botón \textit{Editar} puedes dar de alta nuevos profesores o compartir una misma materia entre varios docentes (el sistema muestra la etiqueta \textit{``Compartida con: ...''} sin sobrescribir a los demás profesores).

\captura{02_datos_profesores.png}{Pestaña Profesores con los 32 docentes oficiales Sep--Dic 2026, categorías FTC/PA/Inglés y materias compartidas.}

Al pulsar el ícono de lápiz (\textit{Editar}) o el enlace \textbf{Configurar Disponibilidad}, se despliega la cuadrícula interactiva de \textbf{Horario de Disponibilidad} (lunes a viernes, de \texttt{07:00} a \texttt{20:45}, respetando los tres recesos institucionales: \texttt{10:40--11:10}, \texttt{13:55--14:00} y \texttt{17:40--18:00}). Las celdas en verde indican bloques habilitados para impartir clase; las celdas blancas bloquean esa franja.

\captura[0.82]{03_datos_disponibilidad.png}{Modal de edición de docente y matriz interactiva de disponibilidad horaria por bloques.}

\textbf{2. Pestaña Materias (Uno o Múltiples Profesores por Materia):} Administra las 68 asignaturas desde \textbf{Inmersión} hasta $10.\textsuperscript{o}$ cuatrimestre (incluyendo \textit{Bases de Datos en la Nube} y \textit{Tutorías}). Cada materia permite asignar \textbf{uno, dos, tres o más profesores} simultáneamente mediante el botón \textbf{+ Profe} o el selector \textbf{Profesores} al agregar/editar la materia; cuando un cuatrimestre abre varios grupos (ej. \texttt{1-1}, \texttt{1-2}, \texttt{1-3} o \texttt{7-1}, \texttt{7-2}), el sistema reparte automáticamente los grupos entre los docentes asignados.

\captura{04_datos_materias.png}{Pestaña Materias mostrando múltiples profesores asignados por asignatura y acceso rápido ``+ Profe''.}

\textbf{3. Pestaña Alumnos (Inmersión y $1.\textsuperscript{o}$ a $10.\textsuperscript{o}$ Cuatrimestre):} Registra la cantidad de estudiantes inscritos en \textbf{Inmersión} y en cada cuatrimestre. Con base en una capacidad estándar de 35 alumnos por aula, el sistema calcula automáticamente cuántos grupos virtuales se deben abrir por cuatrimestre mediante:
\[
\text{Número de grupos del cuatrimestre } q = \left\lceil \frac{\text{Alumnos}_q}{35} \right\rceil
\]
Por ejemplo, 32 alumnos generan 1 grupo (\texttt{1-1}), mientras que 70 alumnos generan 2 grupos (\texttt{7-1} y \texttt{7-2}).

\captura{05_datos_alumnos.png}{Captura de matrícula estudiantil por cuatrimestre para el dimensionamiento automático de grupos.}

\subseccionUPV{Módulo Clásico de Horarios (\texttt{/schedule})}

Este módulo conecta tus datos de \texttt{localStorage} con el endpoint \texttt{POST /api/solve} (motor clásico \texttt{solver.exe} en C++, con respaldo transparente hacia CP-SAT si faltan librerías de compilador).
\begin{enumerate}
    \item Haz clic en el botón azul \textbf{Generar Horario}.
    \item Selecciona el grupo deseado en el menú desplegable \textbf{Seleccionar Grupo}.
    \item Descarga el horario individual o el paquete completo usando los botones \textbf{Excel}, \textbf{PDF}, \textbf{Todos (Excel)} o \textbf{Todos (PDF)}.
\end{enumerate}

\captura{06_horarios_clasico.png}{Horario semanal generado en el módulo clásico con controles de selección de grupo y exportación.}

En la barra central \textbf{Vista}, puedes alternar hacia dos herramientas analíticas adicionales:
\begin{itemize}
    \item \textbf{Matriz:} Muestra la matriz de adyacencia entre cursos coloreando en verde pares que comparten docente, en amarillo pares del mismo grupo y en morado ambos conflictos.
    \item \textbf{Grafo:} Construye un grafo interactivo de fuerzas que enlaza cursos, docentes, grupos y bloques de tiempo para inspeccionar visualmente la densidad de restricciones.
\end{itemize}

\captura{07_horarios_matriz.png}{Vista de Matriz de Adyacencia de Cursos para identificar relaciones de grupo y docente.}

\captura{08_horarios_grafo.png}{Vista de Grafo interactivo de relaciones entre cursos, docentes, grupos y franjas horarias.}

\newpage
% ==============================================================================
% 5. MANUAL DE USUARIO: MÓDULOS DE INTELIGENCIA ARTIFICIAL
% ==============================================================================
\seccionUPV{Manual de Usuario: Módulos de Inteligencia Artificial}

La sección \textbf{Inteligencia Artificial} del menú lateral concentra los motores formales en Python, la importación/exportación en formato Sábana Oficial UPV, el asistente de lenguaje natural y el banco de pruebas comparativo.

\subseccionUPV{Optimizador IA (\texttt{/ia-optimizer})}

Este es el motor principal de generación y validación académica. Su flujo de uso consta de cuatro pasos sencillos:

\begin{enumerate}
    \item \textbf{Elegir la Instancia / Fuente de Datos (R1):}
    \begin{itemize}
        \item \textbf{\texttt{ITIID Sep--Dic 2026 (Excel UPV)}:} Instancia institucional completa con 10 grupos, 25 docentes, franjas fijas de Inglés y bloqueos de posgrado (\texttt{MI}/\texttt{DC}).
        \item \textbf{\texttt{Mis Datos Actuales (LocalStorage)}:} Toma en vivo los profesores, materias y alumnos capturados en el módulo \textbf{Datos}, sin alterar tu almacenamiento original.
        \item \textbf{Instancias de evaluación (\texttt{I1} a \texttt{I4}):} \texttt{I1} (benchmark público ITC-2007), \texttt{I2} (7.\textsuperscript{o} cuatrimestre ITIID con horario manual de referencia), \texttt{I3} (multicuatrimestre con recursos compartidos bloqueados) e \texttt{I4} (escenario de reprogramación).
        \item \textbf{Plantilla Excel / Importar Excel:} Con los botones superiores \textbf{Plantilla Excel} e \textbf{Importar Excel} puedes descargar el formato \texttt{.xlsx} estructurado, llenarlo en Excel y cargarlo directamente con validación en \texttt{pandas}.
    \end{itemize}
    \item \textbf{Seleccionar el Motor de Optimización:}
    \begin{itemize}
        \item \textbf{OR-Tools CP-SAT (R2 --- Motor Principal):} Programación por restricciones con desigualdades booleanas exactas y arranque cálido (\textit{warm-start}).
        \item \textbf{Algoritmo Genético (R3):} Evolución poblacional con selección por torneo, cruce uniforme, mutación y reparación en múltiples corridas independientes (reportando media y desviación estándar).
        \item \textbf{Búsqueda Local (R4):} Construcción voraz con heurística \textbf{MRV} (\textit{Minimum Remaining Values}) seguida de \textbf{Hill Climbing con reinicios aleatorios}.
        \item \textbf{LLM Local (Ollama):} Prueba directa con modelo generativo para contrastar contra los métodos exactos.
    \end{itemize}
    \item \textbf{Elegir el Escenario de Pesos Blandos:}
    \begin{itemize}
        \item \textbf{Escenario A (Prioriza Estudiantes):} Minimiza horas muertas estudiantiles (peso 10) y jornadas mayores a 6 horas diarias, compactando grupos matutinos desde las \texttt{07:00}.
        \item \textbf{Escenario B (Prioriza Docentes):} Minimiza huecos entre clases docentes (peso 10), desbalance de carga y violaciones de preferencias personales (peso 10).
        \item \textbf{Escenario C (Equilibrio):} Balancea de forma simétrica la compacidad estudiantil y las preferencias docentes.
    \end{itemize}
    \item \textbf{Ejecutar y Auditar:} Pulsa \textbf{Ejecutar Motor IA}. Las cuatro tarjetas resumen el dictamen del \textbf{Validador Oficial}: violaciones duras (\texttt{FACTIBLE (0)}), penalización blanda total, huecos de alumnos/docentes y tiempo de cómputo en milisegundos.
\end{enumerate}

\captura{09_ia_optimizador_cpsat.png}{Configuración del Optimizador IA ejecutando OR-Tools CP-SAT sobre I2 con dictamen FACTIBLE (0 violaciones duras).}

En la parte inferior puedes cambiar entre cinco vistas (\textbf{Sábana General}, \textbf{Sábana Docentes}, \textbf{Por Grupo}, \textbf{Por Docente} y \textbf{Por Aula}) y descargar el libro multipestaña con formato institucional mediante \textbf{Excel Completo (R1)} o \textbf{PDF (Sábana)}.

\captura{10_ia_sabana_upv.png}{Vista Sábana General en formato oficial UPV (grupos horizontales L $\cdot$ M $\cdot$ Mi $\cdot$ J $\cdot$ V y franjas de receso).}

\textbf{Reprogramación Estable ante Imprevistos (I4):} Si a mitad del periodo un docente deja de estar disponible un día completo, selecciona el docente y el día bloqueado en la franja \textbf{Reprogramación (I4 --- Estabilidad)} y pulsa \textbf{Reprogramar CP-SAT}. El modelo resuelve nuevamente las restricciones duras penalizando cualquier movimiento innecesario respecto al horario vigente, reportando cuántas sesiones se movieron y el porcentaje del horario que permaneció intacto.

\captura{11_ia_reprogramacion_i4.png}{Resultado de reprogramación mínima I4: bloqueo sobrevenido resuelto con solo 4 sesiones movidas (88.6\% de estabilidad).}

\subseccionUPV{Asistente LLM (\texttt{/ia-llm})}

Este módulo materializa la interacción en lenguaje natural ejecutándose \textbf{100\% en local mediante Ollama} (\texttt{http://127.0.0.1:11434}, modelo predeterminado \texttt{qwen2.5:3b}), sin depender de ninguna API externa ni enviar datos fuera de tu computadora. Aun así, conserva el módulo de anonimización (\texttt{Docente\_01}, \texttt{Aula\_01}) para protección institucional.

\textbf{1. Pestaña R5 --- Extractor de Preferencias y Evaluación F1 (P2):}
\begin{itemize}
    \item Escribe en español las solicitudes de los profesores (por ejemplo, evitar viernes en la tarde, limitar horas diarias o preferir turno matutino) y pulsa \textbf{Extraer Restricciones JSON (R5)}.
    \item Revisa en el panel derecho cada regla traducida, marca o desmarca las casillas de \textbf{Revisión Humana} y haz clic en \textbf{Aprobar y Aplicar al Optimizador} para incorporarlas a la siguiente corrida de CP-SAT, Genético o Búsqueda Local.
    \item Más abajo, el botón \textbf{Evaluar Dataset Anotado (P2)} mide automáticamente \textbf{Precisión}, \textbf{Recall} y \textbf{F1-Score} sobre las 10 oraciones de referencia desglosadas por categoría de restricción.
\end{itemize}

\captura{12_llm_extractor_r5.png}{Pestaña R5 del Asistente LLM: texto libre en español, vista anonimizada y casillas obligatorias de revisión humana.}

\captura{13_llm_evaluacion_f1.png}{Evaluación cuantitativa P2 del extractor R5 mostrando Precisión, Recall y F1 global y por tipo de restricción.}

\textbf{2. Pestaña R6 --- Explicador de Conflictos y Pedidos de Cambio (P4 / I4):}
\begin{itemize}
    \item \textbf{Explicar Último Horario:} Lee el reporte numérico del Validador Común y redacta en español claro un diagnóstico ejecutivo con fortalezas, huecos detectados y recomendaciones de ajuste para la coordinación académica.
    \item \textbf{Pedido de Cambio de Coordinación:} Escribe un imprevisto en lenguaje cotidiano (ej. \textit{``Docente\_04 ya no puede impartir clases los miércoles''}) y pulsa \textbf{Traducir Pedido y Reprogramar con CP-SAT} para traducir el bloqueo y obtener al instante un horario reparado con mínimo desplazamiento.
\end{itemize}

\captura{14_llm_explicador_r6.png}{Pestaña R6: explicación ejecutiva de métricas blandas (izquierda) y reprogramación directa desde texto libre (derecha).}

\textbf{3. Pestaña Prueba P1 --- LLM Directo vs Validador:} Permite demostrar experimentalmente por qué un LLM no debe asignar horarios sin un solver matemático: ejecuta la generación directa con y sin una ronda de retroalimentación del validador, exhibiendo los choques de grupo, aula y recesos que el Validador Común atrapa.

\captura{15_llm_prueba_p1.png}{Pestaña Prueba P1: auditoría empírica de las violaciones duras cometidas por generación directa frente al Validador Común.}

\subseccionUPV{Comparativa y Ética (\texttt{/ia-benchmark})}

Este módulo automatiza la evaluación científica de todos los algoritmos y documenta los controles éticos institucionales:
\begin{itemize}
    \item \textbf{Tabla de Posiciones (P3):} Selecciona la instancia (ej. \texttt{I2}), el escenario de pesos y pulsa \textbf{Ejecutar Comparativa Completa}. El sistema corre los seis enfoques sobre los mismos datos y los clasifica bajo la regla oficial: primero \textbf{cero violaciones duras} y segundo \textbf{menor penalización blanda}. Incluye botón para exportar la tabla comparativa a Excel.
    \item \textbf{Gráficas Comparativas y Síntesis P1--P4:} Grafica en doble eje las violaciones duras frente a la penalización blanda, así como las horas muertas de estudiantes frente a huecos docentes.
    \item \textbf{Pestaña Ética y Checklist:} Presenta la matriz de riesgos éticos, un \textbf{Verificador de Anonimización en vivo} (que muestra cómo se enmascaran tus docentes reales antes de salir a cualquier API) y la lista de verificación interactiva del proyecto.
\end{itemize}

\captura{16_benchmark_leaderboard.png}{Tabla de Posiciones (P3) comparando CP-SAT, Búsqueda Local, Algoritmo Genético, Horario Manual y LLM Directo.}

\captura{17_benchmark_graficas.png}{Gráficas comparativas de violaciones duras, penalización blanda, horas muertas de alumnos y huecos docentes.}

\captura{18_benchmark_etica_anon.png}{Pestaña de gobernanza ética con matriz de mitigación, verificador de anonimización en vivo y checklist del proyecto.}

\newpage
% ==============================================================================
% 6. CÓMO MODIFICAR Y EXTENDER EL SISTEMA PARA AGREGAR MÁS COSAS
% ==============================================================================
\seccionUPV{Cómo Modificar el Sistema para Agregar Más Funciones}

La arquitectura modular del proyecto fue diseñada para que cualquier integrante del equipo o desarrollador futuro pueda ampliar catálogos, reglas de negocio, algoritmos de optimización o pantallas sin romper lo existente. La Tabla~\ref{tab:archivos} resume qué archivo controla cada parte del sistema.

\begin{table}[htbp]
\centering
\small
\caption{Mapa rápido de archivos clave para modificar o extender la aplicación.}
\label{tab:archivos}
\begin{tabularx}{\textwidth}{@{}>{\raggedright\arraybackslash}p{6.2cm}>{\raggedright\arraybackslash}X@{}}
\toprule
\textbf{Archivo del Proyecto} & \textbf{Qué puedes modificar o agregar ahí} \\
\midrule
\texttt{frontend/src/utils/initialData.js} & Catálogo inicial de materias (\texttt{INITIAL\_COURSES}), profesores (\texttt{INITIAL\_TEACHERS}) y bloques horarios (\texttt{TIME\_SLOTS}). \\
\texttt{frontend/src/App.jsx} y \texttt{components/Sidebar.jsx} & Rutas nuevas de navegación (\texttt{<Route>}) y botones del menú lateral izquierdo. \\
\texttt{frontend/src/pages/*.jsx} & Vistas de interfaz de usuario (\texttt{Dashboard}, \texttt{Data}, \texttt{Schedule}, \texttt{IAOptimizer}, \texttt{IALLM}, \texttt{IABenchmark}). \\
\texttt{frontend/src/utils/iaStorage.js} & Configuración por defecto del LLM (\texttt{DEFAULT\_LLM\_CONFIG}), puente \texttt{callIAAction()} y exportación PDF en cliente. \\
\texttt{server/server.js} & Puertos, tiempos máximos de espera (\texttt{PYTHON\_TIMEOUT\_MS}) y comunicación Express $\leftrightarrow$ Python / C++. \\
\texttt{backend\_ia/instances.py} & Instancias oficiales (\texttt{ITIID\_2026}, \texttt{I1}--\texttt{I4}), nuevos periodos cuatrimestrales y anonimizador de nombres. \\
\texttt{backend\_ia/validator.py} & Franjas horarias, pesos de Escenarios (\texttt{SCENARIO\_WEIGHTS}) y reglas de validación duras y blandas. \\
\texttt{backend\_ia/solvers.py} & Modelos de optimización: OR-Tools CP-SAT, Algoritmo Genético y Búsqueda Local (o nuevos metaheurísticos). \\
\texttt{backend\_ia/llm\_service.py} & Prompts, soporte de modelos LLM, reglas NLP de extracción de preferencias y dataset anotado F1. \\
\texttt{backend\_ia/excel\_io.py} & Diseño visual, colores, columnas y pestañas de la Sábana Oficial UPV exportada a Excel (\texttt{.xlsx}). \\
\texttt{backend\_ia/cli.py} & Enrutador principal JSON de acciones Python (\texttt{solve}, \texttt{validate}, \texttt{benchmark}, \texttt{reschedule}, etc.). \\
\texttt{backend/solver.cpp} & Algoritmo clásico de backtracking en C++ (requiere recompilar \texttt{solver.exe} si se edita). \\
\bottomrule
\end{tabularx}
\end{table}

\subseccionUPV{Caso 1: Agregar Nuevas Materias, Docentes, Aulas o un Nuevo Cuatrimestre Oficial}

\begin{itemize}
    \item \textbf{Sin tocar código (vía Excel o Interfaz):} Descarga la plantilla desde \textbf{Optimizador IA $\rightarrow$ Plantilla Excel}, agrega tantas filas de docentes, aulas, materias, bloqueos y preferencias como necesites en las hojas del archivo \texttt{.xlsx}, y cárgala con \textbf{Importar Excel}.
    \item \textbf{Dejar datos fijos en el catálogo inicial web:} Abre \texttt{frontend/src/utils/initialData.js} y añade objetos al arreglo \texttt{INITIAL\_COURSES} o \texttt{INITIAL\_TEACHERS}. Después entra a \textbf{Datos} en el navegador y haz clic en el botón gris \textbf{Predefinidos} para recargar tu catálogo en \texttt{localStorage}.
    \item \textbf{Crear una nueva instancia institucional en Python (ej. \texttt{ITIID\_2027}):}
    \begin{enumerate}
        \item En \texttt{backend\_ia/instances.py}, crea una función \texttt{\_build\_itiid\_ene\_abr\_2027()} retornando el diccionario con \texttt{teachers}, \texttt{rooms}, \texttt{courses}, \texttt{blocked\_slots} y \texttt{preferences}.
        \item Regístrala dentro de \texttt{get\_instance(instance\_id)}:
        \cajaCodigo{if instance\_id == "ITIID\_2027":\\\hspace*{1.2em}return \_build\_itiid\_ene\_abr\_2027()}
        \item En \texttt{frontend/src/pages/IAOptimizer.jsx}, agrega \texttt{<option value="ITIID\_2027">ITIID Ene-Abr 2027</option>} dentro del selector de fuentes de datos.
    \end{enumerate}
\end{itemize}

\subseccionUPV{Caso 2: Agregar una Nueva Restricción Dura o Blanda (o un Escenario D)}

Como el sistema utiliza un \textbf{Validador Común} compartido por todos los motores, toda regla nueva se agrega en solo dos archivos:
\begin{enumerate}
    \item \textbf{Definir pesos o validación en \texttt{backend\_ia/validator.py}:}
    \begin{itemize}
        \item Si deseas crear un \textbf{Escenario D} (por ejemplo, priorizar uso de laboratorios o compactar viernes), agrega la clave \texttt{"D"} al diccionario \texttt{SCENARIO\_WEIGHTS} (líneas 33--61).
        \item Si es una nueva regla dura o métrica blanda, agrega su conteo dentro de \texttt{validate\_schedule()} para que el árbitro oficial la evalúe en cualquier horario.
    \end{itemize}
    \item \textbf{Incorporarla al motor exacto en \texttt{backend\_ia/solvers.py}:}
    \begin{itemize}
        \item Dentro de \texttt{\_build\_and\_solve\_cpsat()}, usa las variables binarias \texttt{x[(cid, slot)]} (que valen \texttt{1} si el curso \texttt{cid} se imparte en el bloque \texttt{slot}).
        \item Para una \textbf{restricción dura} (prohibir algo estrictamente): añade \texttt{model.Add(...)}, por ejemplo \texttt{model.Add(sum(vars\_viernes\_tarde) == 0)}.
        \item Para un \textbf{criterio blando} (penalizar sin prohibir): añade el término ponderado a \texttt{obj\_terms.append(peso * var)}. Los motores Genético y Búsqueda Local heredan automáticamente cualquier cambio hecho en \texttt{validate\_schedule()} porque su función de aptitud consulta al validador.
    \end{itemize}
\end{enumerate}

\subseccionUPV{Caso 3: Agregar un Nuevo Algoritmo de Optimización (ej. Recocido Simulado)}

Si deseas añadir un cuarto algoritmo en Python (por ejemplo, \textit{Simulated Annealing} o \textit{Tabu Search}):
\begin{enumerate}
    \item En \texttt{backend\_ia/solvers.py}, escribe una función \texttt{solve\_annealing(raw\_instance, scenario="C")} que construya una lista de asignaciones \texttt{[\{"course\_id": c, "room\_id": r, "start\_slot": s\}, ...]} y retorne el diccionario estándar:
    \cajaCodigo{%
    return \{\\
    \hspace*{1em}"method": "Recocido Simulado",\\
    \hspace*{1em}"method\_id": "annealing",\\
    \hspace*{1em}"elapsed\_ms": elapsed,\\
    \hspace*{1em}"schedule": best\_schedule,\\
    \hspace*{1em}"validation": validate\_schedule(raw\_instance, best\_schedule, scenario=scenario)\\
    \}%
    }
    \item En \texttt{backend\_ia/cli.py} (sección \texttt{if action == "solve":}), agrega:
    \cajaCodigo{if method == "annealing": return solve\_annealing(instance, scenario=scenario)}
    (Y opcionalmente añádelo a la lista \texttt{results} dentro de \texttt{if action == "benchmark":} para que compita automáticamente en la Tabla de Posiciones).
    \item En \texttt{frontend/src/pages/IAOptimizer.jsx}, agrega \texttt{<option value="annealing">Recocido Simulado</option>} en el selector \textbf{Motor de Optimización}.
\end{enumerate}

\subseccionUPV{Caso 4: Agregar Nuevas Preferencias en Español o Cambiar de Modelo LLM}

\begin{itemize}
    \item \textbf{Usar otro modelo en Ollama (ej. Llama 3.2, Mistral o Gemma):} Descárgalo en terminal con \texttt{ollama pull llama3.2:3b}, entra a \textbf{Asistente LLM}, pulsa \textbf{Detectar modelos locales} y elígelo en el selector. Si quieres dejarlo como predeterminado en el código, cambia \texttt{model: 'qwen2.5:3b'} en \texttt{DEFAULT\_LLM\_CONFIG} dentro de \texttt{frontend/src/utils/iaStorage.js}.
    \item \textbf{Soportar un nuevo tipo de preferencia docente:}
    \begin{enumerate}
        \item En \texttt{backend\_ia/validator.py}, amplía \texttt{preference\_forbidden\_slots(pref)} indicando qué índices de \texttt{slot} penaliza la nueva etiqueta (ej. \texttt{type == "avoid\_friday\_morning"}).
        \item En \texttt{backend\_ia/llm\_service.py}, añade la etiqueta en la instrucción del prompt JSON de \texttt{extract\_preferences()} y su expresión regular de respaldo en el extractor local, además de agregar ejemplos en \texttt{GOLD\_PREFERENCES\_DATASET} para evaluar su F1.
    \end{enumerate}
\end{itemize}

\subseccionUPV{Caso 5: Agregar una Nueva Pantalla / Módulo Web y Acción Backend}

El puente Express (\texttt{/api/ia/action} en \texttt{server/server.js}) reenvía dinámicamente cualquier JSON hacia \texttt{backend\_ia/cli.py}, por lo que agregar un módulo nuevo completo toma cuatro pasos mínimos:
\begin{enumerate}
    \item En \texttt{backend\_ia/cli.py}, agrega un bloque \texttt{if action == "mi\_nueva\_accion": return \{...\}}.
    \item Crea el archivo React \texttt{frontend/src/pages/NuevoModulo.jsx} e invoca tu acción Python con \texttt{await callIAAction(\{ action: 'mi\_nueva\_accion', ... \})}.
    \item En \texttt{frontend/src/App.jsx}, importa \texttt{NuevoModulo} y añade \texttt{<Route path="/nuevo-modulo" element=\{<NuevoModulo />\} />}.
    \item En \texttt{frontend/src/components/Sidebar.jsx}, agrega \texttt{\{ path: '/nuevo-modulo', label: 'Nuevo Módulo', icon: Sparkles \}} al arreglo \texttt{iaNavItems}. Vite recargará la interfaz en caliente (\textit{Hot Module Replacement}) sin reiniciar nada.
\end{enumerate}

\subseccionUPV{Caso 6: Modificar el Solver Clásico en C++ (\texttt{backend/solver.cpp})}

Si editas la heurística de poda o validación en \texttt{backend/solver.cpp}, recompila el ejecutable desde terminal con:
\cajaCodigo{%
g++ backend/solver.cpp -o backend/solver.exe -O2%
}
O bien elimina \texttt{backend/solver.exe} y ejecuta \texttt{run\_project.bat}, el cual detectará la ausencia del binario y lo compilará automáticamente al iniciar.

% ==============================================================================
% 7. SOLUCIÓN DE PROBLEMAS FRECUENTES
% ==============================================================================
\newpage
\seccionUPV{Solución de Problemas Frecuentes (Troubleshooting)}

La Tabla~\ref{tab:errores} concentra los errores más habituales durante la instalación o ejecución y su solución directa.

\begin{table}[htbp]
\centering
\small
\caption{Diagnóstico rápido y solución de incidentes comunes.}
\label{tab:errores}
\begin{tabularx}{\textwidth}{@{}>{\raggedright\arraybackslash}p{4.3cm}>{\raggedright\arraybackslash}p{4.4cm}>{\raggedright\arraybackslash}X@{}}
\toprule
\textbf{Síntoma o Mensaje} & \textbf{Causa Probable} & \textbf{Solución Paso a Paso} \\
\midrule
\texttt{No se pudo iniciar Python} o \texttt{No module named 'ortools'} & Python no está en el \texttt{PATH} de Windows o faltan paquetes \texttt{pip}. & Reinstala Python marcando \textit{Add python.exe to PATH} y ejecuta \texttt{pip install ortools openpyxl pandas requests}, luego reinicia \texttt{server}. \\
\addlinespace[3pt]
\texttt{Failed to fetch} al pulsar \textbf{Ejecutar Motor IA} & El servidor Express (\texttt{server/}) en el puerto \texttt{3001} está apagado. & Abre una terminal en \texttt{server/} y ejecuta \texttt{npm start} (o inicia con \texttt{run\_project.bat}). \\
\addlinespace[3pt]
En \textbf{Horarios} clásico aparece: \textit{``El motor clásico determinó que no hay un horario viable''} (HTTP 422) & Hay varios grupos activos para un mismo docente cuyas horas disponibles en \textbf{Datos} no alcanzan sin traslape. & Amplía los bloques verdes del docente en \textbf{Datos $\rightarrow$ Profesores}, reduce grupos en \textbf{Alumnos}, o resuélvelo en \textbf{Optimizador IA} con \textbf{OR-Tools CP-SAT} (que incluye diagnóstico relajado). \\
\addlinespace[3pt]
\texttt{Ollama Local no detectado en :11434} & Ollama no está instalado o el servicio local está detenido. & Ejecuta \texttt{ollama serve} en terminal, o cambia el selector \textbf{Proveedor LLM} a \textbf{Respaldo Offline} para trabajar sin Ollama. \\
\addlinespace[3pt]
Cambios en \texttt{initialData.js} no aparecen en la tabla de \textbf{Datos} & El navegador sigue leyendo la versión previa guardada en \texttt{localStorage}. & Entra a \textbf{Datos $\rightarrow$ Profesores / Materias} y pulsa el botón gris \textbf{Predefinidos} (o limpia datos del sitio en herramientas de desarrollador \texttt{F12}). \\
\addlinespace[3pt]
Error \texttt{EADDRINUSE: address already in use :::3001} & Ya quedó otra ventana de \texttt{Node Server} abierta en segundo plano. & Cierra la ventana anterior de \texttt{Node Server} antes de volver a abrir \texttt{run\_project.bat}. \\
\bottomrule
\end{tabularx}
\end{table}

\cajaNota{Recomendación final para demostraciones académicas}{Para una presentación ágil ante jurado o docente: (1) abre \textbf{Optimizador IA} y ejecuta \textbf{OR-Tools CP-SAT} sobre \textbf{I2} o \textbf{ITIID Sep--Dic 2026}, mostrando el dictamen \texttt{FACTIBLE (0)} y exportando la \textbf{Sábana General UPV} en Excel; (2) ejecuta una \textbf{Reprogramación I4} en vivo; (3) en \textbf{Asistente LLM} muestra extracción \textbf{R5} con revisión humana, métricas \textbf{F1 (P2)} y la \textbf{Prueba P1}; y (4) cierra en \textbf{Comparativa y Ética} ejecutando la \textbf{Tabla de Posiciones} completa.}

\end{document}
